% Generated draft: requires mathematical and source review.
\documentclass[11pt]{article}
\usepackage[margin=25mm]{geometry}
\usepackage{fontspec}
\setmainfont{Latin Modern Roman}
\usepackage{amsmath,amssymb,mathtools,unicode-math}
\setmathfont{Latin Modern Math}
\usepackage{enumitem,needspace,xurl,hyperref}
\setlength{\parindent}{0pt}
\setlength{\parskip}{5pt}
\setlength{\emergencystretch}{3em}
\newcommand{\N}{\mathbb{N}}
\newcommand{\Z}{\mathbb{Z}}
\newcommand{\Q}{\mathbb{Q}}
\newcommand{\R}{\mathbb{R}}
\newcommand{\C}{\mathbb{C}}
\newcommand{\F}{\mathbb{F}}
\newcommand{\A}{\mathbb{A}}
\newcommand{\PP}{\mathbb{P}}
\newcommand{\E}{\mathbb{E}}
\newcommand{\eps}{\varepsilon}
\newcommand{\dd}{\,\mathrm{d}}
\begin{document}
\section*{UnsolvedMath ID: 30005857}
\textbf{Source problem code:} catalogue entry
\subsection{Definitions and mathematical statement}
For a smooth Riemannian metric \(g\), write \(\operatorname{scal}(g)\) for its scalar curvature. For a real number \(a\), its negative part is \(a_-:=\max\{-a,0\}\). These conventions explain notation in the supplied context; they do not supply missing hypotheses.

The recoverable question is: \emph{``What about rigidity in comparison with sphere?''} The record does not specify the manifold, its dimension, the regularity of the metric under consideration, the comparison sphere or its normalization, the comparison assumptions, or the conclusion meant by rigidity. Consequently, no precise quantified rigidity assertion can be recovered from this record.

The surrounding passage discusses Ricci--DeTurck smoothing and scalar curvature lower bounds. It describes an unspecified system as strictly parabolic and related by diffeomorphisms to Ricci flow in the smooth case. It attributes to Bamler [1] a uniform localization argument giving a desired scalar curvature lower bound for the unique Ricci--DeTurck smoothing as \(t\to0^+\), and to Burkhardt-Guim [3] a generalization to general compact manifolds defining a notion of \(\operatorname{scal}\geq k\) for continuous metrics. These are supplied contextual descriptions, with their hypotheses and reference entries absent.

The passage also attributes to Huang and an unidentified author [6] a smoothing result for a continuous metric \(g_\infty\) that is a \(C^0\) limit of smooth metrics \(g_i\). Its integrability condition contains \((\operatorname{scal}(g_i)-\kappa)_-\), convergence to zero, and displaced fragments \(n/2\) and \(L_{\mathrm{loc}}\); its asserted scalar curvature conclusion contains \(\operatorname{scal}(g_\infty(t))\), \(\kappa\), and \(o(1)\) as \(t\to0\), but the comparison relation is missing. These expressions cannot be restored unambiguously. The passage relates this discussion to Miao's positive mass theorem with corners [12] and suggests that an \(o(1)\) buffer might be removable when \(\kappa\) is constant, using improved approximation from Ricci flow smoothing. This suggestion is preserved as source context, not promoted to a separate conjecture.

The supplied statement-verification note says that the broad sphere-rigidity question introduces a proved rigidity theorem in the official report, and that a remaining precise smoothing question is preserved in the following record. Neither that theorem nor the following record is supplied here. The registry labels this entry \texttt{open}, while marking its statement \texttt{unrecoverable}; these labels do not establish an open mathematical assertion.
\subsection{Short English statement}
This record asks broadly about rigidity in comparison with a sphere, but supplies no hypotheses or intended rigidity conclusion. Its surrounding discussion of scalar curvature and Ricci--DeTurck smoothing is incomplete and mathematically corrupted. The supplied verification note describes the sphere question as introducing a proved theorem, despite the registry's \texttt{open} label. A precise problem statement requires recovery and review of the relevant source passage.
\subsection{Sources}
\begin{itemize}
\item Local UnsolvedMath registry, permanent ID \(30005857\), title \emph{Scalar-Curvature Rigidity Context}, exact source code \texttt{catalogue entry}; collection: Oberwolfach Reports Open Problems. The supplied record gives the proposed year as 2024 and identifies no proposer.
\item Supplied report citation: \emph{Analysis, Geometry and Topology of Positive Scalar Curvature Metrics} (2024), Oberwolfach Reports, pp. 483--566, DOI 10.4171/owr/2024/9. Supplied source URL and literature-source link labelled ``original OWR report'': \url{https://doi.org/10.4171/owr/2024/9}. The extracted passage ends with a report-header fragment and page number 537; its exact relevant passage has not been independently verified.
\item Supplied literature assessment, checked at \texttt{2026-08-22}: the extracted rigidity question does not identify a precise theorem or hypothesis, so no definitive solved or disproved claim can be assigned. This is supplied dated curation, not an independently verified literature assessment; the background explicitly calls for independent verification of time-sensitive status and cited claims.
\item Supplied statement-verification note: the official report uses the broad question to introduce a proved rigidity theorem, and the remaining precise smoothing question is preserved in the following record. This note is supplied provenance, not an independently checked theorem statement.
\item Attribution retained from the excerpt: Bamler [1], Burkhardt-Guim [3], Huang and an unidentified author [6], and Miao [12]. The corresponding bibliographic entries are not supplied, so these reference numbers cannot be resolved here.
\end{itemize}
\end{document}
